\documentclass[12pt]{article}
\usepackage[margin=1in]{geometry}
\usepackage[T1]{fontenc}
\usepackage[utf8]{inputenc}
\usepackage{setspace}
\usepackage{parskip}

\title{The Last Remaining Asymmetry, Part II\\
Why Do Young Women Become More Progressive as Societies Become More Equal?}
\author{Kangbin Yim\thanks{Email: yim@sch.ac.kr (Prof., SCH University, Korea)}}
\date{Jun. 2026}

\begin{document}
\maketitle

For decades, political scientists have attempted to explain why women in many advanced democracies increasingly support progressive political movements. Traditional explanations often focus on policy preferences. Women may favor stronger welfare systems. They may support gender equality initiatives.
They may place greater emphasis on social protection, healthcare, education, or inclusion. These explanations are not necessarily wrong.

But they may not fully explain why this phenomenon appears to become stronger precisely in societies that are already among the most equal, prosperous, and democratic in human history. If equality is increasing, why does the demand for even greater equality continue to intensify? The answer may lie in a paradox hidden inside the success of modern civilization itself.

\section*{The Equality Paradox}

Throughout history, civilization has steadily reduced many forms of asymmetry.
Political rights expanded. Educational opportunities broadened. Economic participation became less dependent upon birth. The distinction between aristocrats and commoners weakened. The distance between rulers and citizens narrowed.

Women gained access to education, employment, property ownership, voting rights, and political participation. In many respects, the history of modern civilization can be understood as a long process of removing structural inequalities. This process has been remarkably successful.

Yet success often reveals what remains. When many inequalities disappear, the remaining ones become more visible. A small stain is difficult to notice on a dirty window. It becomes impossible to ignore on a clean one.

\section*{Why Remaining Differences Become More Visible}

This phenomenon appears repeatedly throughout history. As slavery weakened, the remaining forms of slavery appeared increasingly unacceptable. As racial segregation diminished, the remaining forms of discrimination became more visible. As legal barriers against women disappeared, attention shifted toward subtler forms of inequality. Progress does not eliminate awareness of asymmetry.
It often increases it. The closer a society moves toward equality, the more sensitive people become to the inequalities that remain.

\section*{The Emerging Political Gap}

This observation may help explain an important pattern emerging across many advanced societies. Large international studies have found that women in developed democracies increasingly identify with progressive political positions, while young men often move in the opposite direction. This pattern has been observed across North America, Western Europe, Northern Europe, and parts of East Asia.

Recent comparative studies have shown that the political gap between young men and young women is particularly pronounced in highly developed societies. South Korea has become one of the most frequently cited examples. The question is why.

\section*{The Last Remaining Asymmetry}

One possibility is that the explanation has less to do with political ideology itself and more to do with the changing structure of social competition. Modern societies increasingly operate according to principles of merit, performance, and equal opportunity. Women and men compete in the same universities. They compete in the same job markets. They pursue similar professional careers. They participate in the same political and economic systems.

In other words, the social roles of men and women have become increasingly symmetrical. This represents one of the greatest achievements of modern civilization. Yet one major asymmetry remains. Biology.

Pregnancy occurs within the female body. Childbirth occurs within the female body. Recovery from childbirth occurs within the female body. The physical risks, career interruptions, and long-term biological burdens associated with reproduction remain unequally distributed.

Technology can reduce some of these burdens. Policy can compensate for some of the costs. But the underlying biological asymmetry remains.

\section*{Fairness and Fertility}

This creates a unique tension. The more society succeeds in making everything else equal, the more visible this remaining difference becomes. In a traditional society where many roles are already unequal, biological asymmetry may simply appear as one difference among many.

In a highly equal society, however, it becomes increasingly difficult to ignore.
The question naturally emerges: If men and women are expected to compete equally in education, employment, promotion, income, and social achievement, why does reproduction continue to impose unequal biological costs?

This question is not merely economic. It is not solely about childcare subsidies or housing prices. It is about fairness. And fairness has become one of the defining values of advanced societies.

Viewed from this perspective, declining fertility and increasing progressive identification among young women may not be separate phenomena. They may represent different expressions of the same structural reality. As social equality expands, awareness of remaining asymmetries increases. As awareness increases, demands for further equality increase.

And as reproductive asymmetry becomes more visible, decisions regarding marriage and childbearing become more complicated. What earlier generations often accepted as a natural role may increasingly be experienced as an unequal burden.

\section*{A Question for the Twenty-First Century}

This does not mean that progressive politics causes low fertility. Nor does it mean that fertility decline causes progressive politics. Rather, both may emerge from the same underlying transformation.

A society organized around fairness inevitably becomes more sensitive to the inequalities that survive its own success. Perhaps this is one of the central political questions of the twenty-first century. Not whether societies should pursue equality. But what happens after equality succeeds.

What happens when nearly every social asymmetry has been reduced, yet biological asymmetry remains. What happens when civilization removes most of the visible differences between men and women, but cannot remove pregnancy itself. The answer may shape not only future politics, but also the future of family formation, fertility, and social cohesion.

And it may explain why some of the most equal societies in human history are simultaneously experiencing some of the lowest birth rates ever recorded. The challenge facing advanced societies may no longer be the inequalities of the past. It may be the asymmetries that remain after nearly everything else has been equalized.

\end{document}
